\documentclass[11pt]{article}
\usepackage[margin=1in]{geometry}
\usepackage[T1]{fontenc}
\usepackage{lmodern,amsmath,booktabs}
\usepackage[hidelinks]{hyperref}
\usepackage{xurl}
\setlength{\parindent}{0pt}
\setlength{\parskip}{6pt}
\setlength{\emergencystretch}{3em}
\title{One Perception Path with Predicted Pose Used Only for Initialization}
\author{Pan Dynamics Collision Avoidance Research}
\date{October 3, 2026}
\begin{document}
\maketitle

\section{Implemented decision}
The sole perception path now predicts the target's current relative position
and body heading, initializes a free three-variable LiDAR rectangle fit, and
then supplies the fitted center to the estimator. The old motion-heading
iterations, locked-heading feedback interface, optional pose branch and
algorithm-selection flags have been removed. This decision supersedes the
optional-path statement in the preceding pose-initialization report.

The predicted pose is only a numerical starting value. It contributes no
penalty, fixed orientation, translation bound or frozen choice of a support
axis. In particular, the preceding implementation selected a support axis
from the seed and held it fixed during optimization. That implicit influence
on the objective has been removed. Point-cloud support is now recomputed at
every candidate pose. Position and heading remain free throughout the fit;
declared vehicle length and width remain fixed at 5.0 and 1.9 m in this study.

\section{Point-cloud objective and execution order}
Let $z=(c_x,c_y,\psi)$ and $u_i=R(\psi)^T(p_i-c)$ for observed target points
$p_i$. Define $h=(L/2,W/2)^T$, $e_i=|u_i|-h$, support quantiles $a_j,b_j$
of the current $u_{ij}$, midpoint $m_j=(a_j+b_j)/2$, and coverage
$s_j=(b_j-a_j)/(2h_j)$. With support threshold $\tau=0.8$, the geometry cost is
\[
 \begin{aligned}
 J(z;P)&=\frac1N\sum_i\left[
 \min\{\min_j |e_{ij}|^2,0.25\}
 +20\|\max(e_i,0)\|^2\right]
 +\sum_{j=1}^2 \bigl(w_j(s_j)m_j\bigr)^2,\\
 w_j(s)&=\operatorname{clip}_{[0,1]}\frac{s-\tau}{1-\tau}.
 \end{aligned}
\]
This objective contains the point cloud and candidate pose, but no predicted
pose. The coverage ramp avoids a discontinuous on/off support residual at
the threshold. Its symmetry assumption remains a geometric modeling
assumption, not an estimator prior or an exact vehicle-surface likelihood.

The first fit uses 2/98 percent support quantiles. Points outside that fitted
rectangle by more than 0.15 m are rejected before a second fit with 0/100
percent support quantiles; if too few remain, the first fit is retained.
The outlier selection uses the first fitted geometry. Both stages precede
one estimator update. The existing NumPy damped Gauss--Newton solver has
60 iterations per stage and no bound on total pose correction. A 0.5 m
scaled per-step limit controls numerical steps only.

\path{run_pipeline(args, pose_feedback=adapter)} now requires an adapter with
\path{predict(frame, time)} and \path{update(frame, time, position)}.
The former returns a \path{PosePrior} in current point-cloud ego axes, with
state time equal to the current frame and latest input strictly earlier.
The standalone CLI connects that adapter through
\texttt{--estimator-factory module:callable}; this is dependency wiring,
not a choice of fitting method. Missing or stale priors are rejected.

At first acquisition, no previous target center exists. A current-cloud
face-placement bootstrap supplies the center before the same free pose fit.
The replay retains its existing co-moving initial-heading assumption.
When detection is missing, the estimator receives a missing observation;
the predicted position is never recycled as a synthetic measurement.
Local initialization can choose a solution in an unobserved direction,
but it cannot make partial point support fully observable or guarantee the
global minimum. A rectangle also has a 180-degree body-direction ambiguity.

\section{Validation and observed accuracy}
All 32 current Python behavior tests pass. Coverage includes full default
pipeline wiring, mandatory estimator integration, correction of all pose
variables, nearby seed convergence, underobserved tangential ambiguity,
support recomputation, prefix invariance, stale inputs and missing detections.
Retired heading-only and motion-iteration tests were replaced by the unified
path's behavior checks. No MATLAB source changed and no MATLAB unit tests
were rerun for this change.

Four offline replays reuse frozen actual detector selections and native LiDAR
clusters from the earlier 15 m and 20 m parked-free recordings. The actual
MATLAB NRMM runtime processes each frame. This is not a new YOLO inference
run or a new CARLA capture. The runtime snapshot is commit
\path{26cd6f3998d17fd49f1cd9254245805d17f4d3ae}, isolating these comparisons
from concurrent controller and bound edits. Inputs are the previously seen
Town10HD\_Opt straight following scenario: 300 evaluation frames at 0.02 s,
native GNSS/IMU ego preprocessing and the earlier ideal-compass limitation.
LiDAR/ego-sensor seed pairs were 20261002/20261003 for 15 m and
20261003/20261004 for 20 m. The original hidden parked meshes remain absent.

\begin{center}\small
\begin{tabular}{@{}lrrr@{}}\toprule
 & \multicolumn{3}{c}{RMSE after 2 seconds} \\
Replay & Raw position (m) & Est. position (m) & Est. velocity (m/s) \\\midrule
15 m & 0.4572 & 0.4682 & 0.3561 \\
15 m, 0.3 s blackout & 0.4560 & 0.4658 & 0.3751 \\
20 m & 0.2038 & 0.1394 & 0.3375 \\
20 m, 0.3 s blackout & 0.2044 & 0.1405 & 0.3357 \\
\bottomrule\end{tabular}
\end{center}

The preceding pose fitter gave 0.2648 m and 0.0637 m normal-run position
RMSE after 2 s. Both position results regress when its seed-dependent
support selection is removed. The new implementation satisfies the intended
initialization-only semantics; these measurements do not demonstrate an
accuracy improvement. The change also alters the support objective, so this
comparison cannot attribute the regression to initialization alone.
No parameter was tuned against the truth scores.

All 1200 estimator publications are finite. The interrupted runs preserve
the first 150 measurements and states exactly, have missing detections at
frames 150--164, and reacquire at 165. The supplied priors equal the runtime's
published current pre-update position and heading. Source hashes match the
snapshot recorded before execution; prediction hashes were written before
separate scoring read actor truth. Prediction orchestration rejects attempts
to read the scoring-truth file.

There are 1170 accepted fits. The first stage reports convergence for 1128;
the refinement reports convergence for 1063. The others still return a
finite improved fit, without claiming numerical convergence. Refinement
iteration medians are 9 for 15 m and 7 for 20 m, with maximum 60.
No comparison of cold-start and predicted-start optimization time was run;
the intended acceleration is not yet a measured speedup or a 50 Hz guarantee.
The model and local-solver limitations remain relevant to reducing RMSE.

\section{Artifacts and scope}
The reusable implementation is in
\path{perception/yolo_lidar_target_position.py} and
\path{perception/rectangle_pose_fit.py}. Decision, source provenance,
commands, execution records, audit and scoring summaries accompany this
report in \path{report/ESTIMATOR_POSE_DEFAULT_20261003/}.
Local orchestration, raw predictions and the derived report PDF are in
\path{simulation_output/estimator_pose_default_20261003/}; recorded inputs,
dependencies, model weights and generated binaries are excluded from Git.
The original recordings and prior experiment reports remain preserved.
This finite, previously seen-data replay establishes the documented behavior
only; it supplies no new deterministic observer error certificate.
\end{document}
